\section{Anomaly Detection (Task 2B)}

\subsection{4-Tier ensemble design}
Anomaly detection uses a layered ensemble to capture complementary abnormality types:
\begin{itemize}
	\item \textbf{Z-Score}: global spikes,
	\item \textbf{IQR}: robust outliers,
	\item \textbf{Isolation Forest}: contextual anomalies,
	\item \textbf{LSTM Autoencoder}: temporal anomalies.
\end{itemize}

The design intentionally combines statistical, distributional, contextual, and temporal detectors. Single-method systems often produce unstable alert streams under changing weather and load regimes; ensemble diversity helps maintain alert quality when one detector becomes locally less reliable.

\subsection{Mathematical definitions}
Z-score:
\[
z = \frac{x - \mu}{\sigma}
\]

Autoencoder reconstruction loss:
\[
\mathrm{MSE} = \frac{1}{n}\sum_{i=1}^{n}(x_i - \hat{x}_i)^2
\]

\subsection{Calibration and consensus}
Method outputs were harmonized with a shared target anomaly rate for Isolation Forest contamination and autoencoder fallback quantiles. Autoencoder thresholds use median-plus-MAD first, then quantile fallback when rates move outside the operational band.

This calibration strategy controls alert-volume drift across regions and periods. In practice, operations teams need predictable alert volumes to avoid triage overload; bounded calibration keeps the system actionable while preserving sensitivity to meaningful deviations.

Final escalation uses a high-confidence consensus rule:
\begin{itemize}
	\item \textbf{Anomaly flag = True} when two or more methods agree.
\end{itemize}
This reduces false positives and improves triage.

\subsection{Validation}
Detected anomalies were compared against \texttt{anomaly\_flag}. Consensus alerts improved precision over single-method triggers while keeping useful recall.

Validation was treated as operational rather than purely academic: the objective is to surface incidents worth investigation, not to maximize raw anomaly counts. Consensus rules therefore prioritize higher-confidence incidents for dispatch and root-cause analysis.

For supporting visuals, see the appendix: final 4-method comparison in \autoref{fig:app_anomaly_4method}, training diagnostics in \autoref{fig:app_ae_loss}, and reconstruction diagnostics in \autoref{fig:app_ae_reconstruction}. Aggregated method totals are summarized in \autoref{tab:app_global_method_summary} and \autoref{tab:app_region_method_summary}.